\chapter{Results}\label{ch:results}
% Pure results presentation and interpretation — no design or methodology

\section{Baseline Characterisation (Stage 0)}
\label{sec:results:stage0}
% Tables showing planner performance on original, unmodified domains
% Interpretation: which domains are inherently hard, which planners struggle the most
% Establishes the reference point for all subsequent comparisons

\section{General Prompt Results (Stage 1)}
\label{sec:results:stage1}
% Did a simple "improve this code" prompt work?
% Validation rates, IPC score gains, coverage changes
% Analysis: worked partially, but often generated broken or unhelpful code

\section{Architecture-Aware Prompt Results (Stage 2)}
\label{sec:results:stage2}
% The breakthrough results
% Validation rates, IPC score jumps, coverage improvements
% Cross-testing analysis: proving architectural rules actually matter
% Comparison with Stage 1: interpreting why architecture-aware hints outperform generic ones

\section{Feedback Loop Results (Stage 3)}
\label{sec:results:stage3}
% Validation and convergence analysis
% How the loop caught errors and fixed them iteratively
% Success rate improvement (e.g., ~44\% to ~80\%)
% Iterative learning patterns and failure recovery analysis

\section{Cross-Stage Comparative Analysis}
\label{sec:results:cross_stage}
% Full pipeline progression: Stages 0 → 1 → 2 → 3 side-by-side
% Ablation study isolating the value of each stage
% LLM effectiveness comparison: which model was smartest
% Planner responsiveness: which planner benefited most from good code
% Domain sensitivity: which domains responded best
% Token efficiency and cost analysis

\section{Statistical Validation}
\label{sec:results:statistical_validation}
% Omnibus tests (Friedman): proving overall significance
% Pairwise comparisons (Wilcoxon): Stage 2 vs. Stage 1, Stage 3 vs. Stage 2
% Effect sizes (Cliff's Delta): magnitude of improvements
% Proving to the committee that the results are not coincidence

\section{Comparison with Prior Work}
\label{sec:results:prior_work}
% Direct comparison with Daniel Elis's thesis results
% How the architecture-aware, feedback-loop pipeline outperforms the generic baseline
